\section{Proofs regarding core properties of \refreduction} \label{sec:reduction:proofs}
First we prove that \refreduction can be executed in FPT time parameterized by $|X_C|$.
\begin{lemma} \label{lem:shrink:running:time}
	\Refreduction can be performed in $2^{\Oh(|X_C|)} \cdot n^{\Oh(1)}$ time.
\end{lemma}
\begin{proof}
	Executing \refmark in step~\ref{itm:shrink:mark:b:star} of the algorithm takes $2^{\Oh(|X_C|)} \cdot n^{\Oh(1)}$ time. We proceed by arguing that the remainder of the algorithm can be performed in $n^{\Oh(1)}$ time.
	
	The loop from step~\ref{itm:shrink:loop} forms the computational bottle-neck of the remaining part of the algorithm. It consists of $2 \cdot |X_C \cup B^\ast|^2$ iterations. As $|B^\ast| = |X_C|^{\Oh(1)}$ by construction, this results in $|X_C|^{\Oh(1)}$ iterations. In each of these iterations, we need to check whether $G[X_B \setminus B^\ast]$ connects two given vertices $u$ and $v$ using a path of a given parity $p$. 
	
	This check can be done by iterating over all components $H$ of $G[X_B \setminus B^\ast]$ and returning whether at least one of those connects $u$ and $v$ with a path of parity $p$. Now, it suffices to show that this can be computed for a single arbitrary such component $H$ in polynomial time.
	
	Because every component $H$ of $G[X_B \setminus B^\ast]$ is a connected bipartite subgraph, checking whether an arbitrary such component connects two given vertices $u$ and $v$ with a path of a given parity $p$ is in fact quite simple. If $p = \text{even}$, it suffices to return whether $u$ and $v$ each have a neighbor in the same partite set of $H$, which can clearly be done in polynomial time. To see that this suffices, note that any even path in $H$ starts and ends in the same partite set, since $H$ is bipartite. Moreover, because~$H$ is connected, any two vertices in it are connected by a path.
	
	By similar arguments, if $p = \text{odd}$, it suffices to return whether $u$ and $v$ have neighbors in opposite partite sets of $H$.
\end{proof}
%
Next, we show that \refreduction outputs a strictly smaller graph than its input graph when the OCC it is given is reducible. Recall that we call an OCC $(X_B, X_C, X_R)$ reducible if $|X_B| > g_r(|X_C|)$ for some polynomial $g_r(x)$. The exact definition of this polynomial had been postponed so far. Having reached the statement whose proof relies on this definition, we now specify this exact polynomial.

Its definition is based on \Cref{lem:covering:corollary} in which sets $Z$ and $T$ are specified. Setting the value of $s$ in this lemma to $3$ yields the existence of a constant $c \in \N$ such that $|Z| \leq c \cdot |T|^{8}$ for large enough $|T|$. Given this constant $c$, we define $g_r \colon \N \to \N$ as $g_r(x) = (6(2^8 c+1)^2+2^8 c) \cdot x^{16}$. With this definition, we can prove the following result. 
\begin{lemma} \label{lem:shrink:size}
	Let $G$ be a graph and let $(X_B, X_C, X_R)$ be an OCC of it. If $(X_B, X_C, X_R)$ is reducible, then running \refreduction with these input parameters yields an output graph~$G'$ with $|V(G')| < |V(G)|$.
\end{lemma}
\begin{proof}
	Note that by construction of $G'$, we have $X_C \cup X_R \subseteq V(G) \cap V(G')$, so that we can define $X_B' := V(G') \setminus (X_C \cup X_R)$. Since $X_B = V(G) \setminus (X_C \cup X_R)$, it suffices to show that $|X_B'| < |X_B|$. To do so, we separately bound the sizes of $B^\ast$ and $V(G') \setminus V(G)$ of which $X_B'$ is the union. 
	
	We express bounds in terms of $|X_C|$ and the values of $c$ and $g_r$ as defined above. Bounding the size of $B^\ast$ using these terms is simple: as it was constructed in \refmark via \Cref{lem:covering:corollary}, it has at most $c \cdot (2 |X_C|)^8$ vertices.
	
	To bound the size of $V(G') \setminus V(G)$, note that it includes at most $6$ vertices for every pair of vertices from $X_C \cup B^\ast$. Hence, $|V(G') \setminus V(G)| \leq 6 \cdot |X_C \cup B^\ast|^2 = 6 \cdot (|X_C| + |B^\ast|)^2$. Inserting the bound $|B^\ast| \leq c \cdot (2 |X_C|)^8$ into this expression, we obtain $|V(G') \setminus V(G)| \leq 6 \cdot \left((2^8 c+1) \cdot |X_C|^8 \right)^2 = 6 \cdot (2^8 c+1)^2 \cdot |X_C|^{16}$.
	
	Combining the bounds on $|B^\ast|$ and $|V(G') \setminus v(G)|$, we get $|X_B'| \leq \left(6(2^8 c+1)^2+2^8 c\right) \cdot |X_C|^{16} = g_r(|X_C|)$. As $(X_B, X_C, X_R)$ is a reducible OCC by assumption, it follows that $|X_B| > g_r(|X_C|) \geq |X_B'|$.
\end{proof}
